\documentclass[10pt,a4paper]{article}
\usepackage[top=1.4cm, bottom=1.4cm, left=1.6cm, right=1.6cm, headheight=14pt, footskip=18pt]{geometry}
\usepackage[utf8]{inputenc}
\usepackage{amsmath,amsfonts,amssymb}
\usepackage{graphicx}
\usepackage{tikz}
\usepackage{circuitikz}
\usetikzlibrary{calc,patterns}
\usepackage[most]{tcolorbox}
\usepackage{enumitem}
\usepackage{fancyhdr}
\usepackage{booktabs}
\usepackage{tabularx}
\usepackage{colortbl}
\usepackage{hyperref}

% --- Palet Warna Resmi UNIB ---
\definecolor{unibblue}{RGB}{0, 32, 96}
\definecolor{unibgold}{RGB}{197, 160, 89}
\definecolor{unibgreen}{RGB}{34, 112, 60}
\definecolor{boxbg}{RGB}{248, 250, 253}
\definecolor{linegray}{RGB}{210, 215, 225}

% --- Pengaturan Header & Footer ---
\pagestyle{fancy}
\fancyhf{}
\lhead{\footnotesize\textbf{\color{unibblue}Universitas Bengkulu} \textbar\ Program Studi S1 Teknik Elektro}
\rhead{\footnotesize\color{gray}Persamaan Diferensial (Genap 2026)}
\lfoot{\footnotesize\color{gray}Problem Set OBE \textbar\ \textbf{Video}: \url{https://youtu.be/5GQHKrvaSwg} \textbar\ \href{https://www.youtube.com/playlist?list=PLS5oOZWXZeTw}{\color{unibblue}\textbf{Playlist}}}
\cfoot{\footnotesize\color{gray}Hal. \thepage\ dari 2}
\rfoot{\footnotesize\color{unibblue}\textbf{Minggu 2: Metode Analitik PDB Orde 1}}
\renewcommand{\headrulewidth}{0.6pt}
\renewcommand{\footrulewidth}{0.4pt}

\begin{document}

% ==================== HALAMAN 1 ====================
\noindent\begin{minipage}[c]{0.68\textwidth}%
  {\large\textbf{\color{unibblue}PROBLEM SET (TUGAS MANDIRI TERSTRUKTUR)}}\\[1mm]
  {\normalsize\textbf{Mata Kuliah: Persamaan Diferensial (2 SKS)}}\\[0.5mm]
  {\small\textbf{Topik Minggu 2:} Metode Analitik PDB Orde 1: Separabel \& Eksak (RC)}%
\end{minipage}%
\hfill%
\begin{minipage}[c]{0.30\textwidth}%
  \begin{tcolorbox}[colback=unibblue!5,colframe=unibblue,boxrule=0.6pt,arc=2pt,left=4pt,right=4pt,top=2pt,bottom=2pt]
    \scriptsize
    \textbf{Nama:} \dotfill\\[1.2mm]
    \textbf{NPM:} \dotfill\\[1.2mm]
    \textbf{Kelas/Klp:} \dotfill\\[1.2mm]
    \textbf{Nilai Final:} \hfill \textbf{/ 100}
  \end{tcolorbox}%
\end{minipage}\par

\vspace{1.5mm}

% Kotak Sub-CPMK & Pedoman Tugas Terstruktur
\begin{tcolorbox}[colback=boxbg,colframe=unibgold!90!black,boxrule=0.6pt,arc=2pt,left=5pt,right=5pt,top=3pt,bottom=3pt,title=\textbf{\scriptsize Capaian Pembelajaran (Sub-CPMK 2) \& Pedoman Pengerjaan Tugas}]
  \scriptsize
  \textbf{Sub-CPMK 2:} Mahasiswa mampu mengidentifikasi dan menyelesaikan PDB orde 1 secara analitik (separabel, eksak, dan faktor integrasi) serta memodelkan dinamika pengisian muatan kapasitor RC.\\
  \textbf{Video Perkuliahan (YouTube):} \url{https://youtu.be/5GQHKrvaSwg} \hfill \textbf{Playlist:} \href{https://www.youtube.com/playlist?list=PLS5oOZWXZeTw}{\color{unibblue}\textbf{bit.ly/pd-unib-playlist}}\\\textbf{Ketentuan Tugas:} Kerjakan secara mandiri pada kertas folio/A4 bergaris atau laporan digital rapi. Lampirkan penurunan rumus analitis lengkap dan cetak visualisasi kode program Julia (Soal C6). Kumpulkan pada awal perkuliahan minggu berikutnya.
\end{tcolorbox}

\vspace{1.5mm}

\noindent{\textbf{\color{unibblue}\large BAGIAN I: FONDASI KONSEPTUAL \& PENURUNAN MATEMATIS}}\par
\vspace{1mm}

% Soal C1
\noindent\textbf{1. (C1 - Mengingat \textbar\ Bobot: 10 Poin)}\par\vspace{0.8mm}
{\footnotesize Sebutkan definisi matematis Persamaan Diferensial Terpisahkan (*Separable ODE*) dan tuliskan formulasi uji eksak Euler untuk bentuk diferensial $M(x,y)dx + N(x,y)dy = 0$.}\par

\vspace{2.5mm}

% Soal C2
\noindent\textbf{2. (C2 - Memahami \textbar\ Bobot: 15 Poin)}\par\vspace{0.8mm}
{\footnotesize Tinjau proses pengosongan kapasitor pada rangkaian RC tanpa sumber luar: $R \frac{dq}{dt} + \frac{q}{C} = 0$. Jelaskan secara fisis mengapa laju peluruhan muatan $\frac{dq}{dt}$ berbanding lurus dengan muatan tersisa $q(t)$, dan apa makna fisis konstanta waktu $\tau = RC$.}\par

\vspace{2.5mm}

% Soal C3
\noindent\textbf{3. (C3 - Menerapkan \textbar\ Bobot: 20 Poin)}\par\vspace{0.8mm}
\noindent\begin{minipage}[t]{0.60\textwidth}%
  {\footnotesize Pada rangkaian pengisian kapasitor di samping, kapasitor awalnya kosong $q(0) = 0$. Sakelar ditutup pada $t=0$. Dengan parameter $V_s = 24\text{ V}$, $R = 50\text{ k}\Omega$, dan $C = 20\,\mu\text{F}$, turunkan persamaan diferensial muatan $q(t)$ secara separabel dan tentukan fungsi tegangan kapasitor $v_C(t)$ untuk seluruh $t \ge 0$.}%
\end{minipage}%
\hfill%
\begin{minipage}[t]{0.37\textwidth}%
  \centering
  \resizebox{0.95\linewidth}{!}{%
  \begin{circuitikz}[american, scale=0.68, transform shape]
    \draw (0,0) to[battery2, l=$V_s$, invert] (0,1.8)
          to[switch, l={$t{=}0$}] (1.6,1.8)
          to[R, l=$R$] (3.2,1.8)
          to[C, l=$C$, v=$v_C(t)$] (3.2,0)
          to[short] (0,0);
    \draw[->, >=stealth, thick, unibblue] (1.4,0.8) arc (160:-160:0.35);
    \node at (1.8,0.8) [unibblue, font=\tiny\bfseries] {$i(t)$};
  \end{circuitikz}%
  }%
\end{minipage}\par

\vfill

% Kotak Formula Rujukan Teknis & Standar Industri
\begin{tcolorbox}[colback=unibblue!3,colframe=unibblue,colbacktitle=unibblue,coltitle=white,boxrule=0.6pt,arc=2pt,left=6pt,right=6pt,top=3pt,bottom=3pt,title=\textbf{\scriptsize FORMULA RUJUKAN TEKNIS \& STANDAR INDUSTRI TERKAIT}]
  \scriptsize
  \begin{itemize}[leftmargin=*, nosep]
  \item Hukum Arus Kirchhoff (KAK): $\sum i = 0$ \quad\textbar\quad Karakteristik Kapasitor: $i_C(t) = C \frac{dv_C}{dt} = \frac{dq}{dt}$
  \item Syarat Keeksakan Euler: $\frac{\partial M(x,y)}{\partial y} = \frac{\partial N(x,y)}{\partial x}$ untuk bentuk $M\,dx + N\,dy = 0$
  \item Konstanta Waktu Rangkaian RC: $\tau = R C$ (detik) \quad\textbar\quad Tegangan Kapasitor: $v_C(t) = V_s \left(1 - e^{-t/\tau}\right)$
  \item Energi Elektrostatika Tersimpan: $W_C(t) = \frac{1}{2} C v_C^2(t)$ (Joule)
\end{itemize}
\end{tcolorbox}

\newpage
% ==================== HALAMAN 2 ====================

\noindent{\textbf{\color{unibblue}\large BAGIAN II: ANALISIS SISTEM, EVALUASI STANDAR, \& KOMPUTASI JULIA}}\par
\vspace{1mm}

% Soal C4
\noindent\textbf{4. (C4 - Menganalisis \textbar\ Bobot: 20 Poin)}\par\vspace{0.8mm}
{\footnotesize Analisis Paradoks Efisiensi Energi Pengisian RC: (a) Hitung energi total yang dikeluarkan sumber tegangan $W_{\text{in}} = \int_0^\infty V_s i(t) dt$; (b) Hitung energi elektrostatika yang tersimpan pada kapasitor $W_C = \frac{1}{2} C V_s^2$; (c) Buktikan secara matematis bahwa disipasi kalor pada resistor selalu tepat $W_R = 50\% W_{\text{in}}$, independen terhadap nilai resistansi $R$!}\par

\vspace{3.5mm}

% Soal C5
\noindent\textbf{5. (C5 - Mengevaluasi \textbar\ Bobot: 20 Poin)}\par\vspace{0.8mm}
{\footnotesize Tinjau persamaan medan fluks: $(3x^2 + 2xy + y^2)dx + (x^2 + 2xy + 3y^2)dy = 0$. Evaluasi apakah persamaan diferensial ini memenuhi syarat eksak. Jika eksak, temukan fungsi keluarga kurva integral $\phi(x,y) = C$ secara sistematis langkah demi langkah!}\par

\vspace{3.5mm}

% Soal C6
\noindent\textbf{6. (C6 - Merancang \& Komputasi Julia \textbar\ Bobot: 15 Poin)}\par\vspace{0.8mm}
{\footnotesize Rancang algoritma numerik eksplisit (Metode Euler) dalam bahasa \textbf{Julia} untuk menyelesaikan pengisian kapasitor non-linier: $\frac{dv_C}{dt} = \frac{V_s - v_C}{R \cdot C(v_C)}$ di mana kapasitansi bervariasi terhadap tegangan: $C(v_C) = C_0 (1 + 0{,}1 v_C)$. Tuliskan skrip fungsi pembaruan waktu dan plot hasilnya.}\par

\vfill

% Tabel Rekapitulasi Skor OBE & Lembar Catatan Umpan Balik Dosen
\noindent\begin{minipage}[b]{0.62\textwidth}%
  \centering
  \scriptsize
  \begin{tabular}{|c|c|c|c|c|c|c|}
    \hline
    \rowcolor{unibblue}
    \textbf{\color{white}C1} & \textbf{\color{white}C2} & \textbf{\color{white}C3} & \textbf{\color{white}C4} & \textbf{\color{white}C5} & \textbf{\color{white}C6} & \textbf{\color{white}Total} \\
    \rowcolor{unibblue!10}
    10 Pts & 15 Pts & 20 Pts & 20 Pts & 20 Pts & 15 Pts & 100 Pts \\
    \hline
    & & & & & & \\[3mm]
    \hline
  \end{tabular}%
\end{minipage}%
\hfill%
\begin{minipage}[b]{0.35\textwidth}%
  \centering
  \scriptsize
  \textbf{Verifikasi Dosen / Asisten:}\\[6mm]
  (\dotfill)\\[1mm]
  \tiny Tanggal: \dotfill
\end{minipage}\par

\vspace{2mm}
\begin{tcolorbox}[colback=white,colframe=unibblue!40!black,colbacktitle=unibblue!10,coltitle=unibblue,boxrule=0.5pt,arc=2pt,left=5pt,right=5pt,top=3pt,bottom=3pt,title=\textbf{\tiny Catatan Evaluasi \& Umpan Balik Dosen Pengampu}]
  \tiny
  \textbf{Rubrik Pemeriksaan Pengerjaan Tugas Mandiri:}\par\vspace{0.8mm}
  $\square$ Penurunan matematis langkah-demi-langkah runtut (C1--C4) \hfill $\square$ Validasi analitis \& standar industri IEEE/IEC/SPLN akurat (C5)\\[0.8mm]
  $\square$ Skrip program Julia mandiri \& grafik visualisasi terlampir (C6) \hfill $\square$ Kerapian format laporan \& ketepatan waktu pengumpulan\\[2.0mm]
  \textbf{Catatan / Koreksi Khusus Dosen:}\\[1.6cm]
\end{tcolorbox}

\end{document}
